\documentclass[11pt,a4paper]{article}
\usepackage[margin=25mm]{geometry}
\usepackage{amsmath,amssymb,amsthm,booktabs,longtable}
\newtheorem{lemma}{Lemma}\newtheorem{theorem}[lemma]{Theorem}
\title{Unique overgroups for elements of orders 16 and 5}
\author{Human Compression project}\date{}
\begin{document}\maketitle
Let $G=O'N$. Throughout, a row is a conjugacy class of subgroups in
the fixed complete 581-class roster of Connor and Leemans. We use
that roster and the previously acquired character evaluations as
explicit borrowed inputs. We prove the generation step needed for
four entries of the original map enumeration; three are additional
to the earlier accepted arguments.

\begin{lemma}\label{incidence}
Let $H,K\leq G$, and let $n_H,n_K$ be their conjugacy-class lengths.
If a fixed $K$ contains $m$ subgroups in the $G$-class of $H$, then
the number of conjugates of $K$ containing a fixed $H$ is
$n_Km/n_H$.
\end{lemma}
\begin{proof}
Count incident pairs $(H',K')$ with $H'\leq K'$ in the two fixed
classes, first by $K'$ and then by $H'$. Conjugation makes the number
over each $H'$ constant.
\end{proof}
In particular, a subgroup with exactly one immediate proper
overgroup inherits that overgroup's unique maximal overgroup, if
the latter exists: every maximal overgroup contains an immediate
one. This propagation is justified only when uniqueness has been
proved. Arbitrary cover paths cannot be multiplied as incidence counts.

The roster has exactly two cyclic classes of order 16, rows 510 and
512. The finite certificate in Table~\ref{closure} contains both and
is closed under all minimal-overgroup entries in the roster. Every
overgroup of either cyclic class therefore lies in the displayed
set: take a chain of immediate overgroups and use induction.
Among these rows, the only ones having order divisible by 5 are
\[
1,\quad7,\quad15,\quad16,\quad33,\quad34.
\]
This is an existence test in a class diagram. It uses no products of
cover multiplicities and makes no assertion that distinct paths
represent distinct subgroups. The reader checks the literal row set,
every outgoing minimal-overgroup entry, and the orders directly.

Rows 7, 15 and 16 are maximal in $G$. Rows 33 and 34 have order 2880
and their only minimal-overgroup class is row 7, of order 161280.
The class lengths are respectively $160005384$ and $2857239$, and a
fixed row-7 subgroup contains 56 subgroups in each of rows 33 and 34.
Lemma~\ref{incidence} gives
\[
\frac{2857239\cdot56}{160005384}=1.
\]
Thus a subgroup in either row has one immediate overgroup, which is
already maximal in $G$.

\begin{theorem}
Every proper subgroup of $G$ containing elements of orders 16 and 5
has a unique maximal overgroup.
\end{theorem}
\begin{proof}
The order-16 element places its subgroup in the certified overgroup
set. The order-5 element makes the subgroup order divisible by 5.
For a proper subgroup the resulting possibilities are precisely
rows 7, 15, 16, 33 and 34. Each has a unique maximal overgroup by the
preceding paragraph.
\end{proof}

For the ordered pairs $(r,s)$ satisfying $|r|=p$, $|s|=q$ and
$|rs|=2$, write $x_L(p,q)$ for their number in a group $L$, and
$e_G(p,q)$ for the number generating $G$. If one order is 16 and the
other is divisible by 5, the theorem and the shared subtraction
argument give
\[
 e_G(p,q)=x_G(p,q)-\sum_{M}[G:M]x_M(p,q),
\]
where $M$ runs over all thirteen maximal conjugacy classes. Division
by $|\operatorname{Aut}G|=921631011840$ counts the ordered generating
orbits, with no further duality quotient. The reused exact character
evaluations give:
\begin{center}\begin{tabular}{rrrrr}\toprule
$p$&$q$&$x_G/|\operatorname{Aut}G|$&$(x_G-e_G)/|\operatorname{Aut}G|$&$n_G$\\\midrule
5 & 16 & $1996$ & $30$ & 1966\\
10 & 16 & $17836$ & $14$ & 17822\\
15 & 16 & $15840$ & $32$ & 15808\\
16 & 20 & $35672$ & $28$ & 35644\\
\bottomrule\end{tabular}\end{center}

The intermediate columns are normalized pair counts. All four final
values agree with the original Table~1. The $(15,16)$ entry was already
proved by the maximal-generation criterion; the other three were
unresolved. No new chirality, principal-refutation or conjecture
credit is claimed. No new finite group acquisition is used.

\begin{table}[ht]\centering
\caption{An upward-closed set containing the two cyclic order-16 classes.}\label{closure}
\begin{tabular}{rl}\toprule Order&Rows\\\midrule
460815505920 & 1\\
3753792 & 2, 3\\
1876896 & 4, 5\\
161280 & 7\\
32928 & 9, 10\\
14880 & 15, 16\\
10752 & 21\\
2880 & 33, 34\\
1568 & 53, 54\\
1536 & 55, 56\\
1344 & 65, 66\\
784 & 93, 95\\
768 & 99, 100\\
672 & 115, 116, 117, 118\\
512 & 142\\
384 & 157, 158\\
256 & 203, 204, 205, 206, 207\\
128 & 266, 267, 273, 274, 275, 276\\
64 & 343, 344, 345, 346, 362, 364\\
32 & 426, 427, 434, 435, 452, 453\\
16 & 510, 512\\
\bottomrule\end{tabular}
\end{table}

\begin{thebibliography}{9}
\bibitem{CL} T. Connor and D. Leemans, \emph{Algorithmic enumeration
of regular maps}, Ars Math. Contemp. 10 (2016), 211--222, Table~1.
\bibitem{roster} T. Connor and D. Leemans, \emph{The subgroup lattice
of $O'N$}, fixed 23-page edition. Exact native pages and documentary
bindings accompany the finite certificate.
\end{thebibliography}
\end{document}
